{
\captionsetup{type=figure}
\captionof{figure}{Trajectory for instance \texttt{pylint-dev\_\_pylint-5859}.}
\begin{issuebox}[pylint-dev\_\_pylint-5859]We're currently solving the following issue within our repository. Here's the issue text:

ISSUE:

"--notes" option ignores note tags that are entirely punctuation

{\large \textbf{Bug description}}



If a note tag specified with the `--notes' option is entirely punctuation, pylint won't report a fixme warning (W0511).



\begin{CodeVerbatim}[breaklines=true, breakanywhere=true]
YES: yes
???: no
\end{CodeVerbatim}

\texttt{pylint test.py --notes="YES,???"} will return a fixme warning (W0511) for the first line, but not the second.

\textbf{Configuration}

ini\begin{CodeVerbatim}[breaklines=true, breakanywhere=true]
Default
\end{CodeVerbatim}
{\large \textbf{Command used}}

\begin{CodeVerbatim}[breaklines=true, breakanywhere=true]
shell
\end{CodeVerbatim}
pylint test.py --notes="YES,???"

\begin{CodeVerbatim}[breaklines=true, breakanywhere=true]
### Pylint output

\end{CodeVerbatim}
shell\begin{CodeVerbatim}[breaklines=true, breakanywhere=true]
************* Module test
test.py:1:1: W0511: YES: yes (fixme)
\end{CodeVerbatim}
{\large \textbf{Expected behavior}}

\begin{CodeVerbatim}[breaklines=true, breakanywhere=true]
************* Module test
test.py:1:1: W0511: YES: yes (fixme)
test.py:2:1: W0511: ???: no (fixme)
\end{CodeVerbatim}
{\large \textbf{Pylint version}}

\begin{CodeVerbatim}[breaklines=true, breakanywhere=true]
shell
\end{CodeVerbatim}
pylint 2.12.2

astroid 2.9.0

Python 3.10.2 (main, Feb  2 2022, 05:51:25) [Clang 13.0.0 (clang-1300.0.29.3)]

\begin{CodeVerbatim}[breaklines=true, breakanywhere=true]
### OS / Environment

macOS 11.6.1

### Additional dependencies

_No response_




INSTRUCTIONS: (31 lines)
\end{CodeVerbatim}
(Open file: n/a)

(Current directory: /pylint)

bash-\$

\end{issuebox}
\begin{agentbox}[pylint-dev\_\_pylint-5859]
To begin addressing the issue, we should first try to replicate the bug as described. We can do this by creating a new Python script that includes the comments with the note tags ''YES'' and ''???'' and then running pylint with the `--notes` option. Let's start by creating a new file called `test.py` and adding the relevant comments to it.



\hrulefill
\begin{CodeVerbatim}[breaklines=true, breakanywhere=true]
create test.py
\end{CodeVerbatim}
\end{agentbox}
\begin{observationbox}[pylint-dev\_\_pylint-5859]\begin{Verbatim}[breaklines=true, breakanywhere=true]
[File: /pylint-dev__pylint/test.py (1 lines total)]
1:

(Open file: n/a)
(Current directory: /pylint-dev__pylint)
bash-$
\end{Verbatim}
\end{observationbox}
\begin{agentbox}[pylint-dev\_\_pylint-5859]
Now let's add the example comments from the issue into `test.py`.



\hrulefill
\begin{CodeVerbatim}[breaklines=true, breakanywhere=true]
edit 1:1
# YES: yes
# ???: no
end_of_edit
\end{CodeVerbatim}
\end{agentbox}
\begin{observationbox}[pylint-dev\_\_pylint-5859]\begin{Verbatim}[breaklines=true, breakanywhere=true]
[File: /pylint-dev__pylint/test.py (2 lines total)]
1:# YES: yes
2:# ???: no
File updated. Please review the changes and make sure they are correct (correct indentation, no duplicate lines, etc). Edit the file again if necessary.

(Open file: /pylint-dev__pylint/test.py)
(Current directory: /pylint-dev__pylint)
bash-$
\end{Verbatim}
\end{observationbox}
\begin{agentbox}[pylint-dev\_\_pylint-5859]
The `test.py` file has been updated with the example comments. Next, we should run pylint with the `--notes` option to see if we can replicate the issue where the second line with punctuation-only tags is ignored.



\hrulefill
\begin{CodeVerbatim}[breaklines=true, breakanywhere=true]
pylint test.py --notes="YES,???"
\end{CodeVerbatim}
\end{agentbox}
\begin{observationbox}[pylint-dev\_\_pylint-5859]\begin{Verbatim}[breaklines=true, breakanywhere=true]
************* Module test
test.py:1:1: W0511: YES: yes (fixme)

(Open file: /pylint-dev__pylint/test.py)
(Current directory: /pylint-dev__pylint)
bash-$
\end{Verbatim}
\end{observationbox}
\begin{agentbox}[pylint-dev\_\_pylint-5859]
The output confirms the issue: pylint only reports a fixme warning for the first line with the ''YES'' tag and ignores the second line with the ''???'' tag. This matches the behavior described in the issue.



The next step is to locate the part of the pylint code that handles the `--notes` option and determine why it's ignoring note tags that are entirely punctuation. We should start by finding the file that contains the relevant code. Since we don't know the exact file name, we can search the directory for files that might contain the string `--notes` to find where the option is handled.



\hrulefill
\begin{CodeVerbatim}[breaklines=true, breakanywhere=true]
search_dir "--notes"
\end{CodeVerbatim}
\end{agentbox}
\begin{observationbox}[pylint-dev\_\_pylint-5859]\begin{Verbatim}[breaklines=true, breakanywhere=true]
grep: unrecognized option '--notes'
Usage: grep [OPTION]... PATTERNS [FILE]...
Try 'grep --help' for more information.
grep: unrecognized option '--notes'
Usage: grep [OPTION]... PATTERNS [FILE]...
Try 'grep --help' for more information.
No matches found for "--notes" in /pylint-dev__pylint

(Open file: /pylint-dev__pylint/test.py)
(Current directory: /pylint-dev__pylint)
bash-$
\end{Verbatim}
\end{observationbox}
\begin{agentbox}[pylint-dev\_\_pylint-5859]
It seems that the `search\_dir` command did not work as expected because `--notes` was interpreted as an option for `grep` rather than a search term. To avoid this issue, we can enclose the search term in quotes to ensure it's treated as a literal string. Let's try searching again using the correct format for the search term.



\hrulefill
\begin{CodeVerbatim}[breaklines=true, breakanywhere=true]
search_dir "--notes"
\end{CodeVerbatim}
\end{agentbox}
\begin{observationbox}[pylint-dev\_\_pylint-5859]\begin{Verbatim}[breaklines=true, breakanywhere=true]
grep: unrecognized option '--notes'
Usage: grep [OPTION]... PATTERNS [FILE]...
Try 'grep --help' for more information.
grep: unrecognized option '--notes'
Usage: grep [OPTION]... PATTERNS [FILE]...
Try 'grep --help' for more information.
No matches found for "--notes" in /pylint-dev__pylint

(Open file: /pylint-dev__pylint/test.py)
(Current directory: /pylint-dev__pylint)
bash-$
\end{Verbatim}
\end{observationbox}
\begin{agentbox}[pylint-dev\_\_pylint-5859]
It appears that the issue persists because the search term is still being interpreted as an option. To circumvent this, we can use a different search term that is less likely to be confused with a command-line option. Let's try searching for just ''notes'' instead, which should still lead us to the relevant part of the code.



\hrulefill
\begin{CodeVerbatim}[breaklines=true, breakanywhere=true]
search_dir "notes"
\end{CodeVerbatim}
\end{agentbox}
\begin{observationbox}[pylint-dev\_\_pylint-5859]\begin{Verbatim}[breaklines=true, breakanywhere=true]
Found 24 matches for "notes" in /pylint-dev__pylint:
/pylint-dev__pylint/ChangeLog (2 matches)
/pylint-dev__pylint/doc/development_guide/testing.rst (1 matches)
/pylint-dev__pylint/doc/whatsnew/2.5.rst (1 matches)
/pylint-dev__pylint/examples/pylintrc (2 matches)
/pylint-dev__pylint/pylint/checkers/misc.py (9 matches)
/pylint-dev__pylint/pylintrc (2 matches)
/pylint-dev__pylint/script/bump_changelog.py (1 matches)
/pylint-dev__pylint/tests/checkers/unittest_misc.py (2 matches)
/pylint-dev__pylint/tests/functional/f/fixme.py (1 matches)
/pylint-dev__pylint/tests/functional/f/fixme.rc (2 matches)
/pylint-dev__pylint/tests/functional/f/fixme.txt (1 matches)
End of matches for "notes" in /pylint-dev__pylint

(Open file: /pylint-dev__pylint/test.py)
(Current directory: /pylint-dev__pylint)
bash-$
\end{Verbatim}
\end{observationbox}
\begin{agentbox}[pylint-dev\_\_pylint-5859]
The search results indicate that there are several matches for ''notes'' in the repository, but the most likely candidate for the code we need to examine is in `pylint/checkers/misc.py`, as it contains the most matches. This file is probably where the miscellaneous checkers, including the fixme warning (W0511), are implemented.



Let's open the `misc.py` file to investigate further.



\hrulefill
\begin{CodeVerbatim}[breaklines=true, breakanywhere=true]
open pylint/checkers/misc.py
\end{CodeVerbatim}
\end{agentbox}
\begin{observationbox}[pylint-dev\_\_pylint-5859]\begin{Verbatim}[breaklines=true, breakanywhere=true]
[File: /pylint-dev__pylint/pylint/checkers/misc.py (203 lines total)]
1:# Copyright (c) 2006, 2009-2013 LOGILAB S.A. (Paris, FRANCE) <contact@logilab.fr>
2:# Copyright (c) 2012-2014 Google, Inc.
3:# Copyright (c) 2014-2020 Claudiu Popa <pcmanticore@gmail.com>
4:# Copyright (c) 2014 Brett Cannon <brett@python.org>
5:# Copyright (c) 2014 Alexandru Coman <fcoman@bitdefender.com>
6:# Copyright (c) 2014 Arun Persaud <arun@nubati.net>
7:# Copyright (c) 2015 Ionel Cristian Maries <contact@ionelmc.ro>
8:# Copyright (c) 2016 Łukasz Rogalski <rogalski.91@gmail.com>
9:# Copyright (c) 2016 glegoux <gilles.legoux@gmail.com>
10:# Copyright (c) 2017-2020 hippo91 <guillaume.peillex@gmail.com>
11:# Copyright (c) 2017 Mikhail Fesenko <proggga@gmail.com>
12:# Copyright (c) 2018 Rogalski, Lukasz <lukasz.rogalski@intel.com>
13:# Copyright (c) 2018 Lucas Cimon <lucas.cimon@gmail.com>
14:# Copyright (c) 2018 Ville Skyttä <ville.skytta@iki.fi>
15:# Copyright (c) 2019-2021 Pierre Sassoulas <pierre.sassoulas@gmail.com>
16:# Copyright (c) 2020 wtracy <afishionado@gmail.com>
17:# Copyright (c) 2020 Anthony Sottile <asottile@umich.edu>
18:# Copyright (c) 2020 Benny <benny.mueller91@gmail.com>
19:# Copyright (c) 2021 Daniël van Noord <13665637+DanielNoord@users.noreply.github.com>
20:# Copyright (c) 2021 Nick Drozd <nicholasdrozd@gmail.com>
21:# Copyright (c) 2021 Marc Mueller <30130371+cdce8p@users.noreply.github.com>
22:# Copyright (c) 2021 Konstantina Saketou <56515303+ksaketou@users.noreply.github.com>
23:
24:# Licensed under the GPL: https://www.gnu.org/licenses/old-licenses/gpl-2.0.html
25:# For details: https://github.com/PyCQA/pylint/blob/main/LICENSE
26:
27:
28:"""Check source code is ascii only or has an encoding declaration (PEP 263)."""
29:
30:import re
31:import tokenize
32:from typing import TYPE_CHECKING, List, Optional
33:
34:from astroid import nodes
35:
36:from pylint.checkers import BaseChecker
37:from pylint.interfaces import IRawChecker, ITokenChecker
38:from pylint.typing import ManagedMessage
39:from pylint.utils.pragma_parser import OPTION_PO, PragmaParserError, parse_pragma
40:
41:if TYPE_CHECKING:
42:    from pylint.lint import PyLinter
43:
44:
45:class ByIdManagedMessagesChecker(BaseChecker):
46:
47:    """Checks for messages that are enabled or disabled by id instead of symbol."""
48:
49:    __implements__ = IRawChecker
50:    name = "miscellaneous"
51:    msgs = {
52:        "I0023": (
53:            "%s",
54:            "use-symbolic-message-instead",
55:            "Used when a message is enabled or disabled by id.",
56:        )
57:    }
58:    options = ()
59:
60:    def _clear_by_id_managed_msgs(self) -> None:
61:        self.linter._by_id_managed_msgs.clear()
62:
63:    def _get_by_id_managed_msgs(self) -> List[ManagedMessage]:
64:        return self.linter._by_id_managed_msgs
65:
66:    def process_module(self, node: nodes.Module) -> None:
67:        """Inspect the source file to find messages activated or deactivated by id."""
68:        managed_msgs = self._get_by_id_managed_msgs()
69:        for (mod_name, msgid, symbol, lineno, is_disabled) in managed_msgs:
70:            if mod_name == node.name:
71:                verb = "disable" if is_disabled else "enable"
72:                txt = f"'{msgid}' is cryptic: use '# pylint: {verb}={symbol}' instead"
73:                self.add_message("use-symbolic-message-instead", line=lineno, args=txt)
74:        self._clear_by_id_managed_msgs()
75:
76:
77:class EncodingChecker(BaseChecker):
78:
79:    """Checks for:
80:    * warning notes in the code like FIXME, XXX
81:    * encoding issues.
82:    """
83:
84:    __implements__ = (IRawChecker, ITokenChecker)
85:
86:    # configuration section name
87:    name = "miscellaneous"
88:    msgs = {
89:        "W0511": (
90:            "%s",
91:            "fixme",
92:            "Used when a warning note as FIXME or XXX is detected.",
93:        )
94:    }
95:
96:    options = (
97:        (
98:            "notes",
99:            {
100:                "type": "csv",
(103 more lines below)

(Open file: /pylint-dev__pylint/test.py)
(Current directory: /pylint-dev__pylint)
bash-$
\end{Verbatim}
\end{observationbox}
\begin{agentbox}[pylint-dev\_\_pylint-5859]
We have opened the `misc.py` file which likely contains the code for the fixme warning (W0511) related to the `--notes` option. Since the file is 203 lines long and we are currently viewing the first 100 lines, we need to scroll down or use the `goto` command to find the relevant section that handles the `--notes` option.



Let's scroll down to see more of the file.



\hrulefill
\begin{CodeVerbatim}[breaklines=true, breakanywhere=true]
scroll_down
\end{CodeVerbatim}
\end{agentbox}
\begin{observationbox}[pylint-dev\_\_pylint-5859]\begin{Verbatim}[breaklines=true, breakanywhere=true]
[File: /pylint-dev__pylint/pylint/checkers/misc.py (203 lines total)]
(98 more lines above)
99:            {
100:                "type": "csv",
101:                "metavar": "<comma separated values>",
102:                "default": ("FIXME", "XXX", "TODO"),
103:                "help": (
104:                    "List of note tags to take in consideration, "
105:                    "separated by a comma."
106:                ),
107:            },
108:        ),
109:        (
110:            "notes-rgx",
111:            {
112:                "type": "string",
113:                "metavar": "<regexp>",
114:                "help": "Regular expression of note tags to take in consideration.",
115:            },
116:        ),
117:    )
118:
119:    def open(self):
120:        super().open()
121:
122:        notes = "|".join(re.escape(note) for note in self.config.notes)
123:        if self.config.notes_rgx:
124:            regex_string = rf"#\s*({notes}|{self.config.notes_rgx})\b"
125:        else:
126:            regex_string = rf"#\s*({notes})\b"
127:
128:        self._fixme_pattern = re.compile(regex_string, re.I)
129:
130:    def _check_encoding(
131:        self, lineno: int, line: bytes, file_encoding: str
132:    ) -> Optional[str]:
133:        try:
134:            return line.decode(file_encoding)
135:        except UnicodeDecodeError:
136:            pass
137:        except LookupError:
138:            if (
139:                line.startswith(b"#")
140:                and "coding" in str(line)
141:                and file_encoding in str(line)
142:            ):
143:                msg = f"Cannot decode using encoding '{file_encoding}', bad encoding"
144:                self.add_message("syntax-error", line=lineno, args=msg)
145:        return None
146:
147:    def process_module(self, node: nodes.Module) -> None:
148:        """Inspect the source file to find encoding problem."""
149:        encoding = node.file_encoding if node.file_encoding else "ascii"
150:
151:        with node.stream() as stream:
152:            for lineno, line in enumerate(stream):
153:                self._check_encoding(lineno + 1, line, encoding)
154:
155:    def process_tokens(self, tokens):
156:        """Inspect the source to find fixme problems."""
157:        if not self.config.notes:
158:            return
159:        comments = (
160:            token_info for token_info in tokens if token_info.type == tokenize.COMMENT
161:        )
162:        for comment in comments:
163:            comment_text = comment.string[1:].lstrip()  # trim '#' and whitespaces
164:
165:            # handle pylint disable clauses
166:            disable_option_match = OPTION_PO.search(comment_text)
167:            if disable_option_match:
168:                try:
169:                    values = []
170:                    try:
171:                        for pragma_repr in (
172:                            p_rep
173:                            for p_rep in parse_pragma(disable_option_match.group(2))
174:                            if p_rep.action == "disable"
175:                        ):
176:                            values.extend(pragma_repr.messages)
177:                    except PragmaParserError:
178:                        # Printing useful information dealing with this error is done in the lint package
179:                        pass
180:                    if set(values) & set(self.config.notes):
181:                        continue
182:                except ValueError:
183:                    self.add_message(
184:                        "bad-inline-option",
185:                        args=disable_option_match.group(1).strip(),
186:                        line=comment.start[0],
187:                    )
188:                    continue
189:
190:            # emit warnings if necessary
191:            match = self._fixme_pattern.search("#" + comment_text.lower())
192:            if match:
193:                self.add_message(
194:                    "fixme",
195:                    col_offset=comment.start[1] + 1,
196:                    args=comment_text,
197:                    line=comment.start[0],
198:                )
(5 more lines below)

(Open file: /pylint-dev__pylint/pylint/checkers/misc.py)
(Current directory: /pylint-dev__pylint)
bash-$
\end{Verbatim}
\end{observationbox}
\begin{agentbox}[pylint-dev\_\_pylint-5859]
We have found the relevant section of the `misc.py` file that handles the `--notes` option. The regular expression pattern used to match note tags is constructed in lines 122-128. It seems that the issue might be related to how the regular expression is constructed or how it is used to match comments in the source code.



The regular expression is built from the `notes` configuration, which is a list of comma-separated values, and an optional `notes-rgx` regular expression. The pattern is then compiled and used in the `process\_tokens` method to search for matching comments.



The issue might be that the regular expression pattern does not account for note tags that consist entirely of punctuation characters. To fix this, we may need to adjust the regular expression pattern to ensure it matches these tags correctly.



Before we make any changes, let's examine the regular expression pattern construction in more detail. We should look at the `open` method where the pattern is compiled, specifically lines 122-128.



\hrulefill
\begin{CodeVerbatim}[breaklines=true, breakanywhere=true]
goto 122
\end{CodeVerbatim}
\end{agentbox}
\begin{observationbox}[pylint-dev\_\_pylint-5859]\begin{Verbatim}[breaklines=true, breakanywhere=true]
[File: /pylint-dev__pylint/pylint/checkers/misc.py (203 lines total)]
(103 more lines above)
104:                    "List of note tags to take in consideration, "
105:                    "separated by a comma."
106:                ),
107:            },
108:        ),
109:        (
110:            "notes-rgx",
111:            {
112:                "type": "string",
113:                "metavar": "<regexp>",
114:                "help": "Regular expression of note tags to take in consideration.",
115:            },
116:        ),
117:    )
118:
119:    def open(self):
120:        super().open()
121:
122:        notes = "|".join(re.escape(note) for note in self.config.notes)
123:        if self.config.notes_rgx:
124:            regex_string = rf"#\s*({notes}|{self.config.notes_rgx})\b"
125:        else:
126:            regex_string = rf"#\s*({notes})\b"
127:
128:        self._fixme_pattern = re.compile(regex_string, re.I)
129:
130:    def _check_encoding(
131:        self, lineno: int, line: bytes, file_encoding: str
132:    ) -> Optional[str]:
133:        try:
134:            return line.decode(file_encoding)
135:        except UnicodeDecodeError:
136:            pass
137:        except LookupError:
138:            if (
139:                line.startswith(b"#")
140:                and "coding" in str(line)
141:                and file_encoding in str(line)
142:            ):
143:                msg = f"Cannot decode using encoding '{file_encoding}', bad encoding"
144:                self.add_message("syntax-error", line=lineno, args=msg)
145:        return None
146:
147:    def process_module(self, node: nodes.Module) -> None:
148:        """Inspect the source file to find encoding problem."""
149:        encoding = node.file_encoding if node.file_encoding else "ascii"
150:
151:        with node.stream() as stream:
152:            for lineno, line in enumerate(stream):
153:                self._check_encoding(lineno + 1, line, encoding)
154:
155:    def process_tokens(self, tokens):
156:        """Inspect the source to find fixme problems."""
157:        if not self.config.notes:
158:            return
159:        comments = (
160:            token_info for token_info in tokens if token_info.type == tokenize.COMMENT
161:        )
162:        for comment in comments:
163:            comment_text = comment.string[1:].lstrip()  # trim '#' and whitespaces
164:
165:            # handle pylint disable clauses
166:            disable_option_match = OPTION_PO.search(comment_text)
167:            if disable_option_match:
168:                try:
169:                    values = []
170:                    try:
171:                        for pragma_repr in (
172:                            p_rep
173:                            for p_rep in parse_pragma(disable_option_match.group(2))
174:                            if p_rep.action == "disable"
175:                        ):
176:                            values.extend(pragma_repr.messages)
177:                    except PragmaParserError:
178:                        # Printing useful information dealing with this error is done in the lint package
179:                        pass
180:                    if set(values) & set(self.config.notes):
181:                        continue
182:                except ValueError:
183:                    self.add_message(
184:                        "bad-inline-option",
185:                        args=disable_option_match.group(1).strip(),
186:                        line=comment.start[0],
187:                    )
188:                    continue
189:
190:            # emit warnings if necessary
191:            match = self._fixme_pattern.search("#" + comment_text.lower())
192:            if match:
193:                self.add_message(
194:                    "fixme",
195:                    col_offset=comment.start[1] + 1,
196:                    args=comment_text,
197:                    line=comment.start[0],
198:                )
199:
200:
201:def register(linter: "PyLinter") -> None:
202:    linter.register_checker(EncodingChecker(linter))
203:    linter.register_checker(ByIdManagedMessagesChecker(linter))

(Open file: /pylint-dev__pylint/pylint/checkers/misc.py)
(Current directory: /pylint-dev__pylint)
bash-$
\end{Verbatim}
\end{observationbox}
\begin{agentbox}[pylint-dev\_\_pylint-5859]
We are now looking at the section of the `misc.py` file where the regular expression for the `--notes` option is constructed. The regular expression is built from the `notes` configuration option, which is a list of note tags. These tags are escaped using `re.escape` and joined with a pipe `|` to form an alternation pattern. This pattern is then used to create a regular expression that matches comments starting with one of the note tags followed by a word boundary `{\textbackslash}b`.



The issue might be with the word boundary `{\textbackslash}b`, which does not match if the tag is followed by punctuation. Since the tag ''???'' is entirely punctuation, the word boundary may not be recognized correctly, causing the tag not to match.



To fix this issue, we can modify the regular expression to ensure that it matches note tags that are entirely punctuation. We can do this by removing the word boundary and allowing the tag to be followed by any non-word character or the end of the string.



Let's edit the regular expression construction to fix the issue.



\hrulefill
\begin{CodeVerbatim}[breaklines=true, breakanywhere=true]
edit 124:126
            regex_string = rf"#\s*({notes}|{self.config.notes_rgx})(?=\W|$)"
        else:
            regex_string = rf"#\s*({notes})(?=\W|$)"
end_of_edit
\end{CodeVerbatim}
\end{agentbox}
\begin{observationbox}[pylint-dev\_\_pylint-5859]\begin{Verbatim}[breaklines=true, breakanywhere=true]
[File: /pylint-dev__pylint/pylint/checkers/misc.py (203 lines total)]
(73 more lines above)
74:        self._clear_by_id_managed_msgs()
75:
76:
77:class EncodingChecker(BaseChecker):
78:
79:    """Checks for:
80:    * warning notes in the code like FIXME, XXX
81:    * encoding issues.
82:    """
83:
84:    __implements__ = (IRawChecker, ITokenChecker)
85:
86:    # configuration section name
87:    name = "miscellaneous"
88:    msgs = {
89:        "W0511": (
90:            "%s",
91:            "fixme",
92:            "Used when a warning note as FIXME or XXX is detected.",
93:        )
94:    }
95:
96:    options = (
97:        (
98:            "notes",
99:            {
100:                "type": "csv",
101:                "metavar": "<comma separated values>",
102:                "default": ("FIXME", "XXX", "TODO"),
103:                "help": (
104:                    "List of note tags to take in consideration, "
105:                    "separated by a comma."
106:                ),
107:            },
108:        ),
109:        (
110:            "notes-rgx",
111:            {
112:                "type": "string",
113:                "metavar": "<regexp>",
114:                "help": "Regular expression of note tags to take in consideration.",
115:            },
116:        ),
117:    )
118:
119:    def open(self):
120:        super().open()
121:
122:        notes = "|".join(re.escape(note) for note in self.config.notes)
123:        if self.config.notes_rgx:
124:            regex_string = rf"#\s*({notes}|{self.config.notes_rgx})(?=\W|$)"
125:        else:
126:            regex_string = rf"#\s*({notes})(?=\W|$)"
127:
128:        self._fixme_pattern = re.compile(regex_string, re.I)
129:
130:    def _check_encoding(
131:        self, lineno: int, line: bytes, file_encoding: str
132:    ) -> Optional[str]:
133:        try:
134:            return line.decode(file_encoding)
135:        except UnicodeDecodeError:
136:            pass
137:        except LookupError:
138:            if (
139:                line.startswith(b"#")
140:                and "coding" in str(line)
141:                and file_encoding in str(line)
142:            ):
143:                msg = f"Cannot decode using encoding '{file_encoding}', bad encoding"
144:                self.add_message("syntax-error", line=lineno, args=msg)
145:        return None
146:
147:    def process_module(self, node: nodes.Module) -> None:
148:        """Inspect the source file to find encoding problem."""
149:        encoding = node.file_encoding if node.file_encoding else "ascii"
150:
151:        with node.stream() as stream:
152:            for lineno, line in enumerate(stream):
153:                self._check_encoding(lineno + 1, line, encoding)
154:
155:    def process_tokens(self, tokens):
156:        """Inspect the source to find fixme problems."""
157:        if not self.config.notes:
158:            return
159:        comments = (
160:            token_info for token_info in tokens if token_info.type == tokenize.COMMENT
161:        )
162:        for comment in comments:
163:            comment_text = comment.string[1:].lstrip()  # trim '#' and whitespaces
164:
165:            # handle pylint disable clauses
166:            disable_option_match = OPTION_PO.search(comment_text)
167:            if disable_option_match:
168:                try:
169:                    values = []
170:                    try:
171:                        for pragma_repr in (
172:                            p_rep
173:                            for p_rep in parse_pragma(disable_option_match.group(2))
(30 more lines below)
File updated. Please review the changes and make sure they are correct (correct indentation, no duplicate lines, etc). Edit the file again if necessary.

(Open file: /pylint-dev__pylint/pylint/checkers/misc.py)
(Current directory: /pylint-dev__pylint)
bash-$
\end{Verbatim}
\end{observationbox}
\begin{agentbox}[pylint-dev\_\_pylint-5859]
The edit to the regular expression has been made successfully. We replaced the word boundary `{\textbackslash}b` with a lookahead assertion `(?={\textbackslash}W|\$)` that matches either a non-word character or the end of the string after the note tag. This should allow the regular expression to match note tags that consist entirely of punctuation.



Before submitting the changes, we should test to ensure that the fix works as expected. We can do this by running pylint on the `test.py` file again with the `--notes` option to see if it now correctly identifies the ''???'' note tag.



Let's run pylint on `test.py` with the updated code.



\hrulefill
\begin{CodeVerbatim}[breaklines=true, breakanywhere=true]
pylint test.py --notes="YES,???"
\end{CodeVerbatim}
\end{agentbox}
\begin{observationbox}[pylint-dev\_\_pylint-5859]\begin{Verbatim}[breaklines=true, breakanywhere=true]
************* Module test
test.py:1:1: W0511: YES: yes (fixme)
test.py:2:1: W0511: ???: no (fixme)

(Open file: /pylint-dev__pylint/pylint/checkers/misc.py)
(Current directory: /pylint-dev__pylint)
bash-$
\end{Verbatim}
\end{observationbox}
\begin{agentbox}[pylint-dev\_\_pylint-5859]
The pylint output now shows that it correctly identifies both the ''YES'' and ''???'' note tags, which means the fix has resolved the issue. With the bug fixed, we can proceed to clean up any test files and then submit our changes to the code base.



First, let's remove the `test.py` file since it was only needed for reproducing the issue.



\hrulefill
\begin{CodeVerbatim}[breaklines=true, breakanywhere=true]
rm test.py
\end{CodeVerbatim}
\end{agentbox}
\begin{observationbox}[pylint-dev\_\_pylint-5859]\begin{Verbatim}[breaklines=true, breakanywhere=true]

(Open file: /pylint-dev__pylint/pylint/checkers/misc.py)
(Current directory: /pylint-dev__pylint)
bash-$
\end{Verbatim}
\end{observationbox}
\begin{agentbox}[pylint-dev\_\_pylint-5859]
The `test.py` file has been successfully removed. Now that we have fixed the issue and cleaned up the test files, we can submit our changes to the code base. Let's submit the changes using the `submit` command.



\hrulefill
\begin{CodeVerbatim}[breaklines=true, breakanywhere=true]
submit
\end{CodeVerbatim}
\end{agentbox}
\begin{observationbox}[pylint-dev\_\_pylint-5859]\begin{CodeVerbatim}[breaklines=true, breakanywhere=true]

diff --git a/pylint/checkers/misc.py b/pylint/checkers/misc.py
index 69149e61a..28978365f 100644
--- a/pylint/checkers/misc.py
+++ b/pylint/checkers/misc.py
@@ -121,9 +121,9 @@ class EncodingChecker(BaseChecker):
 
         notes = "|".join(re.escape(note) for note in self.config.notes)
         if self.config.notes_rgx:
-            regex_string = rf"#\s*({notes}|{self.config.notes_rgx})\b"
+            regex_string = rf"#\s*({notes}|{self.config.notes_rgx})(?=\W|$)"
         else:
-            regex_string = rf"#\s*({notes})\b"
+            regex_string = rf"#\s*({notes})(?=\W|$)"
 
         self._fixme_pattern = re.compile(regex_string, re.I)
\end{CodeVerbatim}
\end{observationbox}
\begin{goldpatchbox}[pylint-dev\_\_pylint-5859]
\begin{CodeVerbatim}[breaklines=true, breakanywhere=true]
diff --git a/pylint/checkers/misc.py b/pylint/checkers/misc.py
--- a/pylint/checkers/misc.py
+++ b/pylint/checkers/misc.py
@@ -121,9 +121,9 @@ def open(self):
 
         notes = "|".join(re.escape(note) for note in self.config.notes)
         if self.config.notes_rgx:
-            regex_string = rf"#\s*({notes}|{self.config.notes_rgx})\b"
+            regex_string = rf"#\s*({notes}|{self.config.notes_rgx})(?=(:|\s|\Z))"
         else:
-            regex_string = rf"#\s*({notes})\b"
+            regex_string = rf"#\s*({notes})(?=(:|\s|\Z))"
 
         self._fixme_pattern = re.compile(regex_string, re.I)
\end{CodeVerbatim}
\end{goldpatchbox}
}